\documentclass[11pt,letterpaper]{article}
\usepackage[T1]{fontenc}
\usepackage[utf8]{inputenc}
\usepackage{lmodern}
\usepackage[margin=1in,headheight=15pt]{geometry}
\usepackage{amsmath,amssymb,amsthm,mathtools}
\usepackage{microtype,booktabs,array,longtable,enumitem}
\usepackage[dvipsnames]{xcolor}
\usepackage{fancyhdr}
\usepackage[colorlinks=true,linkcolor=MidnightBlue,citecolor=MidnightBlue,urlcolor=MidnightBlue]{hyperref}
\usepackage[nameinlink,noabbrev]{cleveref}
\usepackage{listings}
\definecolor{Ink}{HTML}{19354A}
\definecolor{Soft}{HTML}{F0F4F6}
\pagestyle{fancy}
\fancyhf{}
\fancyhead[L]{\small\textsc{Surcomplex Analysis}}
\fancyhead[R]{\small Hahn-supported holomorphic functions}
\fancyfoot[C]{\thepage}
\setlength{\parskip}{4pt}
\setlength{\parindent}{1.3em}
\setlist{nosep,leftmargin=*}
\allowdisplaybreaks[2]
\numberwithin{equation}{section}
\newtheorem{theorem}{Theorem}[section]
\newtheorem{proposition}[theorem]{Proposition}
\newtheorem{lemma}[theorem]{Lemma}
\newtheorem{corollary}[theorem]{Corollary}
\theoremstyle{definition}
\newtheorem{definition}[theorem]{Definition}
\newtheorem{example}[theorem]{Example}
\newtheorem{construction}[theorem]{Construction}
\theoremstyle{remark}
\newtheorem{remark}[theorem]{Remark}
\newcommand{\No}{\mathbf{No}}
\newcommand{\K}{\mathbb{K}}
\newcommand{\C}{\mathbb{C}}
\newcommand{\R}{\mathbb{R}}
\newcommand{\N}{\mathbb{N}}
\newcommand{\Z}{\mathbb{Z}}
\newcommand{\OO}{\mathcal{O}}
\newcommand{\mm}{\mathfrak{m}}
\newcommand{\HH}{\mathcal{H}}
\newcommand{\CC}{\mathcal{C}}
\newcommand{\dd}{\,\mathrm{d}}
\newcommand{\st}{\operatorname{st}}
\newcommand{\supp}{\operatorname{supp}}
\newcommand{\Res}{\operatorname{Res}}
\newcommand{\ord}{\operatorname{ord}}
\newcommand{\Hol}{\operatorname{Hol}}
\newcommand{\id}{\operatorname{id}}
\newcommand{\fin}{\operatorname{fin}}
\newcommand{\val}{\operatorname{v}}
\newcommand{\Exp}{\operatorname{Exp}}
\newcommand{\Log}{\operatorname{Log}}
\newcommand{\cint}{\mathop{\int}\nolimits^{\mathrm H}}
\newcommand{\cooint}{\mathop{\oint}\nolimits^{\mathrm H}}
\newcommand{\Tay}{\operatorname{Tay}}
\newcommand{\lead}{\operatorname{lc}}
\newcommand{\smallo}{\mathrm{o}}
\newcommand{\bigo}{\mathrm{O}}
\newcommand{\mono}[1]{t^{#1}}
\hypersetup{pdftitle={Surcomplex Analysis via Hahn-Supported Holomorphic Functions},pdfauthor={Research exposition prepared with OpenAI},pdfsubject={Surreal complex numbers, Hahn summation, preparation, residues, exponential functions}}
\begin{document}
\begin{titlepage}
\vspace*{0.45in}
{\color{Ink}\Large\textsc{A constructive analytic framework}}\par
\vspace{0.4in}
{\color{Ink}\Huge\bfseries Surcomplex Analysis}\par
\vspace{0.18in}
{\LARGE via Hahn-Supported\\[4pt] Holomorphic Functions}\par
\vspace{0.4in}
{\large Taylor calculus, Weierstrass preparation, contour formulas,\\[3pt]
zero counting, and global exponential phenomena}\par
\vspace{0.5in}
\noindent\rule{\textwidth}{0.6pt}
\vspace{0.12in}
\begin{abstract}
\noindent We construct a coherent theory of analytic functions with values in the
surcomplex field $\K=\No[i]$. The basic objects are set-supported Hahn series
with ordinary holomorphic coefficient functions. Taylor evaluation turns these
objects into genuine functions at finite surcomplex arguments, while affine
charts move the construction to arbitrary centers and scales. We prove a strong
Taylor theorem, identity, open-mapping, and maximum-modulus principles, a constructive
Weierstrass preparation and division theorem, a perturbative inverse theorem,
Cauchy and Morera theorems, and a residue theorem whose residues are taken at
actual surcomplex poles. An infinitesimal Rouch\'e theorem counts all zeros,
including nonstandard splittings of multiple zeros. We also construct a
surjective exponential on the finite-imaginary strip, exhibit explicit global
exponential extensions, and prove that any conjugation-compatible global
extension must acquire nonstandard periods.

The assumptions are essential. The full surreal topology admits no
non-eventually-constant convergent set-indexed nets. Ordinary contour paths
therefore cannot supply the proposed integrals. We distinguish Hahn summation
from topological convergence and prove counterexamples to unrestricted
identity, Liouville, and Schwarz principles. The resulting theory is a
specified analytic category, not a claim that all conceivable surcomplex
functions satisfy classical complex analysis.
\end{abstract}
\vfill
\noindent\textbf{Date:} September 21, 2026\par
\noindent\textbf{Status:} Proof-based research exposition and proposed framework.
Classical and published foundations are cited; no priority claim is made for
standard Hahn-series mechanisms or their consequences.
\end{titlepage}
\tableofcontents
\newpage

\section{Purpose, scope, and relationship to existing work}\label{sec:introduction}

A surcomplex number is a formal pair $x+iy$, where $x,y\in\No$ and $i^2=-1$.
The arithmetic is exactly the arithmetic of a quadratic field extension.
Since $\No$ is real closed, $\No[i]$ is algebraically closed. The difficulty is
not algebraic closure: it is deciding what ``analytic'', ``sum'', ``contour'',
and ``continuation'' should mean over a proper-class, non-Archimedean field.
Conway, Gonshor, and Alling provide the basic algebraic and analytic foundations
\cite{Conway,Gonshor,Alling}.

Three structures will be kept separate throughout this article. First, the
surreal-valued modulus gives a very fine local topology and a meaningful
$\varepsilon$--$\delta$ derivative. Second, Conway normal forms give a
support-based notion of infinite summation. Third, ordinary complex domains
provide a geometry on which coefficients can be continued and integrated.
Our central construction combines the second and third structures; the first
then verifies that the functions obtained really are differentiable. The
ordinary complex-analytic inputs applied to coefficients may be found in
Ahlfors \cite{Ahlfors}.

The theory of surreal power-series evaluation is not new. Berarducci and
Mantova explicitly develop summability and local surreal analyticity, including
power series with surreal coefficients \cite[Section 3]{BMgerms}. Their work
on derivations and transseries concerns a further structure on surreal
\emph{numbers} \cite{BMderivations}. We will not identify that derivation with
the derivative of a function of an independent surcomplex variable. The
operator methods of van der Hoeven provide a broad setting for support-based
implicit equations \cite{vdH}. Our preparation proof below gives a direct
specialized recursion, rather than appealing to a completion or to the
convergence of Newton iterates.

Other analytic approaches have different aims. Rubinstein-Salzedo and
Swaminathan study surreal limits and integration using additional constructions
\cite{RSS}. Costin and Ehrlich develop extension and integration methods for a
large class of resurgent functions \cite{CE}. Mantova's January 2026 preprint
proves monotonicity and a Taylor approximation theorem for transseries
\cite{Mantova2026}. None of those results is being replaced here. Also, the
term ``Hahn-holomorphic'' already occurs in M\"uller and Strohmaier's theory of
convergent generalized expansions on complex logarithmic covers \cite{MS}.
To avoid conflation, our objects are called \emph{Hahn-supported holomorphic
functions}: their monomials are surcomplex constants, not functions of the
complex coordinate.

\subsection{The main package}

For an ordinary complex domain $U$, write $U^\sharp$ for its infinitesimal
thickening in the finite surcomplex numbers. The main conclusions are as
follows.

\begin{theorem}[Summary of the construction]\label{thm:summary}
Let $U$ be a connected ordinary complex domain. There is a natural algebra
$\HH(U)$ of set-supported series
\[
 F(z)=\sum_{\gamma\in S} f_\gamma(z)t^\gamma,
 \qquad f_\gamma\in\Hol(U),
\]
with an injective realization as functions $U^\sharp\to\K$. In this algebra:
\begin{enumerate}[label=(\roman*)]
\item derivatives are coefficientwise, and Taylor expansion is an exact Hahn
sum throughout each infinitesimal monad;
\item a nonzero function has isolated zeros in the fine topology, and a
function constant on a nonempty fine-open set is constant on all of
$U^\sharp$;
\item a zero of multiplicity $m$ in the leading coefficient deforms into
exactly $m$ surcomplex zeros in its monad, counted with multiplicity;
\item coefficientwise contours satisfy Cauchy's formula, Morera's theorem,
the residue theorem, and an argument principle counting actual surcomplex
zeros and poles;
\item biholomorphic leading terms admit Hahn-supported inverse perturbations;
\item a genuine local maximum of the surcomplex modulus forces constancy,
but boundedness alone does not imply a Liouville theorem.
\end{enumerate}
All these assertions are proved below with their precise domain hypotheses.
\end{theorem}

The strongest structural step is the preparation theorem in
\cref{sec:preparation}. The residue theorem in \cref{sec:residues} depends on
it: coefficientwise residues at an ordinary center are identified with sums
of residues at the surcomplex poles lying infinitesimally near that center.

\subsection{What is and is not asserted}

The article supplies a rigorous, nontrivial analytic category and works out its
consequences. It does not assert a universal, uniquely determined meaning of
surcomplex analyticity. In particular, the contour integral is a specified
coefficientwise integral, not a Riemann integral along a fine-topologically
continuous real path. Nor does the construction give a unique global
continuation from $\C$ to all of $\K$. Indeed, the explicit exponential
extensions in \cref{sec:globalexp} demonstrate why uniqueness fails.

\section{The surcomplex field, valuation, and size conventions}\label{sec:field}

\subsection{Proper classes and normal forms}

We work in a class theory such as NBG with global choice. All supports,
coefficient collections, recursions, and summation index collections appearing
in an individual construction are sets. The symbols $\No$ and $\K$ denote
proper classes with the stated field operations. We never form a sum indexed
by all surreal numbers.

Every surcomplex number has a unique normal form
\begin{equation}\label{eq:normal}
 z=\sum_{\gamma\in S}c_\gamma t^\gamma,
 \qquad c_\gamma\in\C\setminus\{0\},\quad
 t^\gamma:=\omega^{-\gamma},
\end{equation}
where $S\subseteq\No$ is a well-ordered set. The notation $\omega^a$ here is
Conway's normal-form monomial map. In particular, the notation $t^\gamma$ is
\emph{defined} by this map; it must not silently be reinterpreted as
$\exp(\gamma\log t)$ for an unrelated notion of surreal exponentiation.
We have $t^{\alpha+\beta}=t^\alpha t^\beta$.
The complex normal form follows by combining the real normal forms of the
real and imaginary parts \cite{Conway,Alling}.

For $z\ne0$, set
\[
 \val(z)=\min S,\qquad \lead(z)=c_{\val(z)},\qquad \val(0)=+\infty.
\]
Here $+\infty$ is a valuation symbol, not a surreal number. Products and sums
satisfy
\begin{equation}\label{eq:valuation}
 \val(zw)=\val(z)+\val(w),\qquad
 \val(z+w)\ge\min\{\val(z),\val(w)\}.
\end{equation}
If the two values in the minimum differ, equality holds. We write
$z\prec w$ when $w\ne0$ and $\val(z)>\val(w)$; equivalently,
$|z|<r|w|$ for every positive real $r$.

Any set of exponents is contained in a set-sized ordered additive subgroup
$\Gamma\subseteq\No$. Thus a fixed calculation takes place in a Hahn field
$\C((t^\Gamma))$, enlarged when necessary to contain the supports of new
arguments or finitely many algebraic roots. This is a size reduction, not an
assertion about the topology induced from all of $\K$; that distinction will
matter in \cref{sec:topology}.

\subsection{Conjugation and modulus}

Define
\[
 \overline{x+iy}=x-iy,\qquad
 |z|=(z\overline z)^{1/2}=(x^2+y^2)^{1/2}\in\No_{\ge0}.
\]
The square root is the nonnegative one in the real-closed field $\No$.

\begin{proposition}\label{prop:modulus}
The modulus is multiplicative and satisfies the triangle inequality. Moreover,
$\val(|z|)=\val(z)$ for $z\ne0$, and the leading coefficient of $|z|$ is
$|\lead(z)|$.
\end{proposition}
\begin{proof}
Multiplicativity follows by squaring. For the triangle inequality, the
identity
\[
 (x^2+y^2)(u^2+v^2)-(xu+yv)^2=(xv-yu)^2\ge0
\]
proves Cauchy--Schwarz over any ordered field. Expanding $|z+w|^2$ then gives
$|z+w|^2\le(|z|+|w|)^2$. Positivity allows square roots. Finally, if
$z=ct^\gamma(1+\epsilon)$ with $\epsilon\prec1$, then
$z\overline z=|c|^2t^{2\gamma}(1+\epsilon)(1+\overline\epsilon)$.
Its positive square root has the claimed leading term.
\end{proof}

\subsection{Finite elements, standard parts, and monads}

Let
\[
 \OO=\{z\in\K:|z|<n\text{ for some }n\in\N\},\qquad
 \mm=\{z\in\K:|z|<1/n\text{ for every }n\ge1\}.
\]
Thus $\OO=\{\val\ge0\}$ and $\mm=\{\val>0\}$, with the usual convention at
zero. The constant coefficient gives a ring homomorphism
\[
 \st:\OO\longrightarrow\C,
 \qquad \ker\st=\mm.
\]
Every finite surcomplex number is uniquely $a+\epsilon$, with $a\in\C$ and
$\epsilon\in\mm$. The monad of $a$ is $a+\mm$. For $U\subseteq\C$, define
\begin{equation}\label{eq:halo}
 U^\sharp=\st^{-1}(U)=\{a+\epsilon:a\in U,\ \epsilon\in\mm\}.
\end{equation}

\section{Why ordinary topology does not supply complex analysis}\label{sec:topology}

\subsection{The fine topology}

The fine topology has balls $|z-a|<r$, where $r$ ranges over all positive
surreals. Equivalently, its neighborhood base can be given by
$\val(z-a)>\gamma$, with $\gamma\in\No$. Indeed, valuation inequalities
imply sufficiently weak modulus inequalities, and conversely a sufficiently
small monomial radius gives any prescribed valuation inequality.
Valuation balls are clopen by \eqref{eq:valuation}. We use this
class-indexed basis and its neighborhood predicates; no collection of all
open subclasses is required as an object of class theory.

\begin{lemma}[Set-sized scales have a positive lower bound]\label{lem:setlower}
If $A\subseteq\No_{>0}$ is a set, there is $d\in\No_{>0}$ with $d<a$ for
every $a\in A$.
\end{lemma}
\begin{proof}
The surreal cut $\{0\mid A\}$ is legitimate because its left and right sides
are sets and $0<a$ for every $a\in A$. Its value is strictly between them.
\end{proof}

\begin{theorem}[No nontrivial set-indexed convergent nets]\label{thm:nets}
Every set-sized subset of $\K$ is uniformly discrete for the fine topology.
Every convergent net indexed by a set is eventually equal to its limit.
In particular, a convergent ordinary sequence is eventually constant.
\end{theorem}
\begin{proof}
For a set $E\subseteq\K$, apply \cref{lem:setlower} to the set of nonzero
distances $|x-y|$ with $x,y\in E$. Balls with the resulting radius separate
all distinct elements of $E$.

If $(x_j)_{j\in J}$ converges to $x$ and $J$ is a set, apply the same lemma to
$\{|x_j-x|:x_j\ne x\}$. A ball about $x$ with the resulting radius contains
none of the $x_j\ne x$, so convergence forces eventual equality.
\end{proof}

\begin{corollary}\label{cor:paths}
Every fine-continuous map from an ordinary connected real interval to $\K$
is constant. The usual inclusion $\C\hookrightarrow\K$ is not continuous
when $\C$ has its Euclidean topology.
\end{corollary}
\begin{proof}
The image of the interval is a set, hence a discrete subspace. A continuous
image of a connected space cannot contain two points in a discrete space.
For the inclusion, the intersection of an infinitesimal ball about a complex
number with $\C$ is a singleton, not an ordinary open neighborhood.
\end{proof}

This also explains a subtle point about set-sized Hahn subfields. Their
\emph{intrinsic} Hahn valuation topology can be nondiscrete, but their
subspace topology inherited from the full class $\K$ is discrete. An external
surreal radius can be smaller than every nonzero element of the whole
set-sized field.

Each monad is fine-open and fine-closed. So is every halo $U^\sharp$, because
its membership is locally constant on the finite elements and the infinite
elements also have neighborhoods disjoint from it. Consequently the classical
curve $\partial U$ used later is a \emph{shadow contour}, not the fine
topological boundary of $U^\sharp$.

\subsection{Derivatives are nevertheless meaningful}

For a function on a fine-open class, write $F'(z)=L$ when
\begin{equation}\label{eq:derivative}
 \forall e\in\No_{>0}\ \exists d\in\No_{>0}\quad
 0<|h|<d\ \Longrightarrow\
 \left|\frac{F(z+h)-F(z)}h-L\right|<e.
\end{equation}
This is a genuine all-scales condition. It cannot be replaced by testing
ordinary sequences. Algebraic manipulations prove the usual sum, product,
and chain rules once the corresponding derivatives exist.

\begin{example}[Differentiability alone is too weak]\label{ex:locallyconstant}
Define $L:\K\to\K$ by $L(z)=0$ for finite $z$ and $L(z)=1$ for infinite $z$.
Both regions are fine-open, so $L$ is locally constant. It is everywhere
complex differentiable, indeed locally represented by a constant power series,
and $L'=0$. Nevertheless it is not constant. It is also bounded by $1$.
Thus neither an unrestricted identity theorem nor an unrestricted Liouville
theorem follows from fine differentiability, even with local power-series
representations.
\end{example}

\section{Hahn summation and infinitesimal power series}\label{sec:sums}

\subsection{Summability is a support condition}

\begin{definition}\label{def:summable}
A set-indexed family $(z_j)_{j\in J}$ of surcomplex numbers is Hahn summable
if the union of their supports is well ordered and each exponent occurs in
only finitely many members of the family. Its sum is obtained by adding the
finitely many contributions to each coefficient.
\end{definition}

This convention permits arbitrary regrouping of an already jointly summable
family. The existence of two iterated sums is not a substitute for joint
summability. We use the standard support facts behind Hahn series and
Neumann's lemma \cite{Neumann,vdH,BMgerms}.

\begin{lemma}[Support calculus]\label{lem:neumann}
Let $A,B$ be well-ordered subsets of an ordered abelian group.
\begin{enumerate}[label=(\roman*)]
\item $A+B$ is well ordered; a fixed element has only finitely many
decompositions $a+b$ with $a\in A$, $b\in B$.
\item If $B\subseteq\Gamma_{>0}$, its additive monoid
$M(B)=\{0\}\cup\bigcup_{n\ge1}(B+\cdots+B)$ is well ordered. Each element
has only finitely many representations as a finite nondecreasing word in $B$.
\item Consequently $A+M(B)$ is well ordered, and substitution of series
supported in $B$ into a formal power series with complex coefficients is
Hahn summable whenever its substituted constant term is zero.
\end{enumerate}
\end{lemma}
\begin{proof}[Explanation of the support mechanism]
For (i), infinitely many distinct decompositions of the same element would
contain an increasing subsequence of first components and hence a decreasing
subsequence of second components, contradicting well-ordering. The same
subsequence argument proves well-ordering of the sum set. Assertion (ii) is
Neumann's positive-support lemma: positivity prevents insertion of zero-cost
letters, and well-ordering prevents an infinite descending sequence or
infinitely many words with a fixed sum. We use this classical combinatorial
lemma as an input, in its usual finite-word form. Applying (i) to $A$ and
$M(B)$ gives (iii). For a fixed exponent, only finitely many words and hence
only finitely many Taylor coefficients contribute.
\end{proof}

\begin{example}[An exact sum that does not converge topologically]
With $t=\omega^{-1}$,
\[
 \sum_{n\ge0}t^n=\frac1{1-t}
\]
is a Hahn sum. Its ordinary partial sums do not converge in the fine topology:
the remainder is $t^{N+1}/(1-t)$, which is larger than $t^\omega$ for every
finite $N$. The identity is algebraically exact despite that failure.
\end{example}

\subsection{Every complex formal series can be evaluated infinitesimally}

\begin{proposition}\label{prop:formal}
If $P(X)=\sum_{n\ge0}a_nX^n\in\C[[X]]$ and $\epsilon\in\mm$, then
$P(\epsilon)$ is a well-defined Hahn sum. Substitution preserves addition,
multiplication, and formal composition when the inner series has zero
constant coefficient. For fixed $\epsilon\ne0$, the map
$P\mapsto P(\epsilon)$ is injective.
\end{proposition}
\begin{proof}
Take $B=\supp(\epsilon)\subseteq\No_{>0}$ in \cref{lem:neumann}.
All expanded terms of $P(\epsilon)$ have support in $M(B)$, and only finitely
many contribute to each exponent. The same argument justifies the joint
families in products and compositions. If $a_m$ is the first nonzero
coefficient, then
\[
 P(\epsilon)=a_m\epsilon^m(1+\eta),\qquad \eta\in\mm,
\]
because every remaining term has value at least $(m+1)\val(\epsilon)$.
This is nonzero, proving injectivity.
\end{proof}

No ordinary radius of convergence is needed. For example,
$\sum_{n\ge0}n!\epsilon^n$ exists for every infinitesimal $\epsilon$, even
though the corresponding complex power series has radius zero. It defines a
local infinitesimal function, not an ordinary holomorphic germ. This is one
reason to distinguish formal local analyticity from the coherent coefficient
category introduced next.

\section{Hahn-supported holomorphic functions}\label{sec:category}

\subsection{The algebra and its actual evaluation}

\begin{definition}\label{def:HH}
For an ordinary open set $U\subseteq\C$, let $\HH(U)$ consist of expressions
\begin{equation}\label{eq:HH}
 F=\sum_{\gamma\in S} f_\gamma t^\gamma,
 \qquad f_\gamma\in\Hol(U),
\end{equation}
where $S\subseteq\No$ is a well-ordered set. Zero coefficient functions may
be omitted. Addition is coefficientwise, and multiplication is Hahn
convolution. An expression is \emph{positive-supported} if $S\subseteq
\No_{>0}$, and \emph{nonnegative-supported} if $S\subseteq\No_{\ge0}$.
For nonzero $F$, its global value $v_U(F)$ is the least exponent with a
nonzero coefficient function, not necessarily the value of $F$ at every point.
\end{definition}

A single support is required on all of $U$. This global support requirement
is a deliberate coherence hypothesis. Restrictions are allowed, but we do not
claim arbitrary local expressions glue to one globally supported expression;
an arbitrary union of local supports need not be well ordered.

\begin{theorem}[Taylor realization]\label{thm:realization}
For $F$ as in \eqref{eq:HH}, the formula
\begin{equation}\label{eq:evaluation}
 \widehat F(a+\epsilon)=
 \sum_{\gamma\in S}\sum_{n\ge0}
 \frac{f_\gamma^{(n)}(a)}{n!}\,t^\gamma\epsilon^n,
 \qquad a\in U,\quad\epsilon\in\mm,
\end{equation}
is jointly Hahn summable. It gives an injective algebra homomorphism from
$\HH(U)$ into the class of functions $U^\sharp\to\K$.
\end{theorem}
\begin{proof}
The expanded support is contained in $S+M(\supp\epsilon)$.
The support lemma proves well-ordering and finite multiplicity. For sums and
products, Taylor's product identities and joint summability identify
coefficientwise multiplication with multiplication of the values.
At a standard point $a$, the normal form is simply
$\widehat F(a)=\sum_\gamma f_\gamma(a)t^\gamma$. If the realized function
vanishes everywhere, uniqueness of normal forms gives $f_\gamma(a)=0$ for
every $a\in U$ and every $\gamma$. Thus each coefficient function is zero.
\end{proof}

We henceforth write $F(z)$ for $\widehat F(z)$. A scalar $c\in\K$ belongs to
$\HH(U)$ as its normal form with constant coefficient functions.

\begin{proposition}[Units and reduction]\label{prop:units}
Suppose
\begin{equation}\label{eq:normalized}
 F=t^s(f_0+R),\qquad R\text{ positive-supported},
\end{equation}
and $f_0$ has no zeros on $U$. Then $F$ is nowhere zero on $U^\sharp$ and
\begin{equation}\label{eq:inversegeometric}
 F^{-1}=t^{-s}f_0^{-1}\sum_{n\ge0}(-R/f_0)^n\in\HH(U).
\end{equation}
For a nonnegative-supported $G$, $\st(G(a+\epsilon))=g_0(a)$, where $g_0$
is its exponent-zero coefficient.
\end{proposition}
\begin{proof}
The geometric expansion is summable by the positive-support lemma. At a
point of shadow $a$, the leading term of $t^{-s}F$ is $f_0(a)\ne0$.
For the reduction statement, every Taylor term with positive exponent or
positive Taylor degree is infinitesimal; the only surviving constant
coefficient is $g_0(a)$.
\end{proof}

\subsection{Classical functions embed canonically}

For $f\in\Hol(U)$, let $f^\sharp$ be the realization with only the
exponent-zero coefficient $f$. This is the canonical infinitesimal Taylor
extension. It preserves algebraic identities and, wherever ordinary
composition is defined,
\[
 (f\circ g)^\sharp=f^\sharp\circ g^\sharp.
\]
Indeed, both sides are the same formal Taylor substitution in each monad.
Its standard part is $f\circ\st$.

For instance,
\[
 e^{a+\epsilon}=e^a\sum_{n\ge0}\frac{\epsilon^n}{n!},\qquad
 \sin(a+\epsilon)=\sum_{n\ge0}\frac{\sin^{(n)}(a)}{n!}\epsilon^n
\]
are unambiguous for $a\in\C$, $\epsilon\in\mm$. On a chosen classical
logarithm domain, the same construction extends the chosen logarithm branch.
A divergent formal power series on a monad need not arise this way.

\section{Differentiation, Taylor expansion, and rigidity}\label{sec:calculus}

\subsection{The strong Taylor theorem}

Define coefficientwise derivatives by
$D^kF=\sum_\gamma f_\gamma^{(k)}t^\gamma$.

\begin{theorem}[Strong Taylor calculus]\label{thm:taylor}
Let $F\in\HH(U)$, $z=a+\epsilon\in U^\sharp$, and $h\in\mm$. Then
\begin{equation}\label{eq:strongtaylor}
 F(z+h)=\sum_{n\ge0}\frac{(D^nF)(z)}{n!}h^n
\end{equation}
is a jointly summable identity after full expansion. The fine derivative of
$F$ is $DF$, and all higher derivatives are $D^kF$.
If $s=\min S$, the remainder after degree $N$ satisfies
\begin{equation}\label{eq:remainder}
 \val\left(F(z+h)-\sum_{n=0}^{N}\frac{(D^nF)(z)}{n!}h^n\right)
 \ge s+(N+1)\val(h).
\end{equation}
\end{theorem}
\begin{proof}
Expand $f_\gamma$ formally at $a$ and replace its variable by
$\epsilon+h$. The binomial formula gives the two-variable Taylor identity.
Every fully expanded term has support in
\[
 S+M(\supp\epsilon\cup\supp h),
\]
and the positive-support lemma proves joint summability. Hence regrouping by
the power of $h$ is legitimate and yields \eqref{eq:strongtaylor}.
Each term in the degree-$N$ remainder contains at least $N+1$ factors of $h$,
while every other nonconstant factor has positive value. This proves
\eqref{eq:remainder}.

For $N=1$, the difference quotient minus $(DF)(z)$ has value at least
$s+\val(h)$. Given a positive surreal error bound $e$, choose a positive
surreal $\beta$ with $s+\beta>\val(e)$. For $0<|h|<t^\beta$, we have
$\val(h)\ge\beta$, and the error has value strictly larger than $\val(e)$,
hence modulus smaller than $e$. This proves the fine derivative. Apply the
same argument to $D^kF$ to obtain all orders.
\end{proof}

\begin{remark}[Constants really are constants]
In this theorem, $D(t^\gamma)=0$: the monomial is a coefficient, independent
of $z$. This is not the Berarducci--Mantova derivation on the field of surreal
numbers, in which particular infinite numbers can have nonzero derivative
\cite{BMderivations}. Confusing these operators would invalidate the
coefficient calculus.
\end{remark}

\subsection{Cauchy--Riemann equations and harmonicity}

\begin{proposition}\label{prop:CR}
Suppose $F=u+iv$ is real Fr\'echet differentiable at a point of $\K=\No^2$,
with an $\No$-linear differential and a remainder $\smallo(|h|)$ in the fine
sense. Then it is complex differentiable there if and only if
\[
 u_x=v_y,\qquad u_y=-v_x.
\]
Every $F\in\HH(U)$ satisfies these equations and has
$u_{xx}+u_{yy}=v_{xx}+v_{yy}=0$.
\end{proposition}
\begin{proof}
A real $2\times2$ matrix is multiplication by a surcomplex number precisely
when it has the form
$\left(\begin{smallmatrix}a&-b\\b&a\end{smallmatrix}\right)$.
That is exactly the displayed condition. The remainder estimate then gives
the complex derivative. For our functions, \cref{thm:taylor} supplies the
real differential and commuting higher derivatives, and differentiating the
Cauchy--Riemann equations gives harmonicity.
\end{proof}
The real Fr\'echet hypothesis is important; mere existence of partial
derivatives is not being used as a sufficient condition.

\subsection{Identity theorems and isolated zeros}

\begin{theorem}[Three forms of identity]\label{thm:identity}
Assume $U$ is connected and $F\in\HH(U)$.
\begin{enumerate}[label=(\roman*)]
\item If $F$ vanishes on a subset of $U\subset\C$ having an ordinary complex
accumulation point in $U$, then $F=0$.
\item If all fine derivatives of $F$ vanish at one point $z\in U^\sharp$,
then $F=0$.
\item If $F$ vanishes on a nonempty fine-open subset of $U^\sharp$, then
$F=0$. If $F\ne0$, each of its zeros is fine-isolated.
\end{enumerate}
\end{theorem}
\begin{proof}
For (i), uniqueness of normal forms gives vanishing of every $f_\gamma$ on
the same standard set. The ordinary identity theorem applies coefficientwise.

For (ii), write $z=a+\epsilon$. Apply the strong Taylor formula to $D^kF$,
centered at $z$ with increment $-\epsilon$. All terms vanish, so
$(D^kF)(a)=0$ for every $k$. Uniqueness of normal forms implies
$f_\gamma^{(k)}(a)=0$ for every $\gamma,k$. Ordinary analyticity and
connectedness give $f_\gamma=0$ on $U$.

Vanishing on a fine-open set implies vanishing of all derivatives at any of
its points, giving the first statement of (iii). For the last, let $F(z)=0$
and $F\ne0$. By (ii), there is a least $m\ge1$ with
$A_m=(D^mF)(z)/m!\ne0$. Then
\[
 F(z+h)=h^m(A_m+hB(h)),\qquad \val(B(h))\ge s,
\]
for infinitesimal $h$, by the remainder estimate. Choose $\val(h)$ so large
that $s+\val(h)>\val(A_m)$. The parenthesis is nonzero, so no sufficiently
nearby point other than $z$ is a zero.
\end{proof}

The ordinary accumulation in (i) and the fine isolation in (iii) refer to
different topologies. In particular, $1/n\to0$ in $\C$ but not in $\K$ with
its fine topology. A fine accumulation-point version of the identity theorem
follows immediately from (iii).

\section{Weierstrass preparation and division}\label{sec:preparation}

This section is the main nonlinear construction. Its input is an ordinary
zero of the leading coefficient. Its output is a polynomial over $\K$ whose
roots are exactly the nearby surcomplex zeros.

\subsection{An elementary analytic splitting}

Let $D$ be an ordinary disc centered at $a$, and put $X=z-a$. For a
holomorphic $h$ on $D$, define
\[
 T_m h=\sum_{j=0}^{m-1}\frac{h^{(j)}(a)}{j!}X^j,
 \qquad Q_mh=\frac{h-T_mh}{X^m}.
\]
The quotient has a removable singularity at $a$, and
\begin{equation}\label{eq:splitting}
 h=T_mh+X^mQ_mh
\end{equation}
is the unique decomposition into a polynomial of degree below $m$ and a
multiple of $X^m$ with holomorphic quotient. For $m=0$, preparation merely
asserts that a nonvanishing leading coefficient gives a unit; below $m\ge1$.

\begin{theorem}[Hahn--Weierstrass preparation]\label{thm:preparation}
Suppose $F=t^s(f_0+R)\in\HH(U)$, where $R$ is positive-supported and $f_0$
has a zero of order $m\ge1$ at $a\in U$. Choose a disc $D\Subset U$ on which
\[
 f_0(z)=X^m u_0(z),\qquad u_0(z)\ne0.
\]
There are unique expressions
\begin{align}
 P(X)&=X^m+b_{m-1}X^{m-1}+\cdots+b_0,
       &&b_j\in\mm,\label{eq:prepP}\\
 V(z)&=1+v(z),&&v\text{ positive-supported in }\HH(D),
\end{align}
such that
\begin{equation}\label{eq:prep}
 F(z)=t^s u_0(z)V(z)P(z-a).
\end{equation}
Their positive supports may be chosen in the additive monoid generated by
$\supp(R)$. The factor $V$ is a unit in $\HH(D)$.
\end{theorem}
\begin{proof}
Divide by $t^s u_0$ and write
\[
 t^{-s}F/u_0=X^m+E(z),\qquad E=R/u_0.
\]
Let $S$ be the positive support of $E$ and $M=M(S)$. Seek
\[
 P=X^m+p,\quad p=\sum_{\gamma\in M\setminus\{0\}}p_\gamma(X)t^\gamma,
 \qquad V=1+\sum_{\gamma\in M\setminus\{0\}}v_\gamma(z)t^\gamma,
\]
where each $p_\gamma$ has degree below $m$. The desired equation is
\begin{equation}\label{eq:preprecursionbase}
 E=p+X^m v+pv.
\end{equation}
Proceed by transfinite recursion over the well-ordered set
$M\setminus\{0\}$. At exponent $\gamma$, all product terms use smaller
positive exponents. Set
\begin{equation}\label{eq:preprecursion}
 H_\gamma=E_\gamma-
 \sum_{\substack{\alpha+\beta=\gamma\\\alpha,\beta>0}}
 p_\alpha v_\beta,
 \qquad p_\gamma=T_mH_\gamma,\quad v_\gamma=Q_mH_\gamma.
\end{equation}
The sum is finite. Both operations in \eqref{eq:splitting} preserve
holomorphicity, so the coefficients are holomorphic on the same disc $D$.
No decreasing sequence of radii occurs in this recursion.

The Hahn series formed from these coefficients are legitimate, since their
supports are subsets of $M$. At each exponent, \eqref{eq:preprecursionbase}
holds by construction; hence the full identity holds. Each coefficient
$b_j=\sum_\gamma [X^j]p_\gamma\,t^\gamma$ is infinitesimal. The leading
coefficient of $V$ is $1$, so \cref{prop:units} makes it a unit.

For uniqueness, take two factorizations and the least exponent at which their
$p$ or $v$ coefficients differ, using the union of their well-ordered
positive supports. At that exponent, all smaller product contributions
cancel, leaving $\Delta p+X^m\Delta v=0$. The uniqueness of
\eqref{eq:splitting} forces both differences to vanish, a contradiction.
\end{proof}

\begin{remark}[Why this is not a convergence argument]
The powers of a positive-support perturbation may have leading values
$n\delta$ which are not cofinal in the whole surreal value group. Therefore
neither an ordinary infinite product nor a sequence of Newton approximations
need converge in the fine topology. The proof instead solves one coefficient
at a time on a well-ordered support, with finitely many contributions at each
step. This is the decisive distinction.
\end{remark}

\begin{theorem}[Hahn--Weierstrass division]\label{thm:division}
Let $P=X^m+\sum_{j<m}b_jX^j$ have infinitesimal coefficients $b_j$. For
$H\in\HH(D)$ there are unique $Q\in\HH(D)$ and $A\in\K[X]$ with
$\deg A<m$ such that
\begin{equation}\label{eq:division}
 H=QP+A.
\end{equation}
\end{theorem}
\begin{proof}
Multiplying $H$ by a monomial allows us to assume its support is nonnegative.
Let $M$ be the monoid generated by the positive supports of $H$ and the
finitely many $b_j$. At exponent zero, divide $H_0$ by $X^m$ using
\eqref{eq:splitting}. At a positive exponent $\gamma$, subtract the already
known coefficient of $Q(P-X^m)$ and again use the same splitting. All
nonleading coefficients of $P$ have positive value, so the unknown
coefficient $Q_\gamma$ occurs only in $X^mQ_\gamma$; the recursion is
triangular and finite at each exponent. This constructs $Q$ and $A$.
Undo the monomial scaling. Uniqueness follows from the least-exponent
argument and the uniqueness of ordinary division by $X^m$.
\end{proof}

\subsection{Exact zero lifting}

\begin{theorem}[Multiplicity is conserved in a monad]\label{thm:zerolifting}
Under the assumptions of \cref{thm:preparation}, $F$ has exactly $m$ zeros in
$a+\mm$, counted with their fine analytic multiplicities. They are the roots
$a+\eta_j$ of $P$, and every $\eta_j$ is infinitesimal.
\end{theorem}
\begin{proof}
The algebraic closure of $\K$ factors $P$ into $m$ linear factors, counted
with multiplicity. A root $\eta$ cannot be infinite: the value of $\eta^m$
is then strictly smaller than the value of every $b_j\eta^j$, so cancellation
is impossible. A finite root with nonzero standard part is also impossible,
because reduction of $P(\eta)=0$ would give $\st(\eta)^m=0$.
Thus all roots lie in $\mm$.

The remaining factor $t^s u_0V$ is nonzero throughout $D^\sharp$. Therefore
its product with $P$ has exactly the same zeros. At a root, the strong Taylor
expansion of the unit starts with a nonzero constant. Multiplication cannot
change the first nonzero power of the local coordinate, so analytic and
polynomial multiplicities agree.
\end{proof}

This statement is much stronger than merely lifting simple zeros: a multiple
ordinary zero is allowed to split into several genuinely nonstandard zeros.
The proof works for arbitrary well-ordered positive supports, not only for
integer powers of one infinitesimal parameter.

\section{Inverse and implicit function theorems}\label{sec:inverse}

\subsection{A constructive perturbative inverse}

\begin{theorem}[Biholomorphic leading term]\label{thm:inverse}
Let $f_0:U\to V$ be an ordinary biholomorphism, with inverse $g_0$. Suppose
$F=f_0+R\in\HH(U)$ and $R$ is positive-supported. There is a unique
\[
 G=g_0+S\in\HH(V),\qquad S\text{ positive-supported},
\]
such that $F\circ G=\id$ and $G\circ F=\id$ on the corresponding halos.
In particular, $F:U^\sharp\to V^\sharp$ is a bijection, and
\[
 G'(w)=\frac1{F'(G(w))}.
\]
\end{theorem}
\begin{proof}
Seek $G=g_0+S$ with positive support in the monoid generated by $\supp(R)$.
Substitute into
\begin{equation}\label{eq:inverserecursion}
 f_0(g_0+S)+R(g_0+S)=w.
\end{equation}
Taylor substitution is supportwise legitimate. At positive exponent
$\gamma$, the unknown coefficient $S_\gamma(w)$ occurs linearly as
$f_0'(g_0(w))S_\gamma(w)$. All other contributions involve smaller positive
exponents and there are only finitely many of them. Because
$f_0'\circ g_0$ is nonzero, division yields a holomorphic coefficient on $V$.
Transfinite recursion therefore produces $G$ and proves $F(G(w))=w$.

Reduction gives $\st(F(z))=f_0(\st z)$, so $F$ maps $U^\sharp$ into
$V^\sharp$. If $F(z_1)=F(z_2)$, the shadows agree after applying $f_0$, hence
$z_j=a+\epsilon_j$ for one $a$. Taylor factorization of the difference gives
\[
 F(a+\epsilon_1)-F(a+\epsilon_2)
 = (\epsilon_1-\epsilon_2)\bigl(f_0'(a)+\eta\bigr),
 \qquad \eta\in\mm.
\]
The second factor is a unit. Thus $z_1=z_2$, proving injectivity and the
second inverse identity. The derivative formula follows from the chain rule.
The same difference argument gives uniqueness of the inverse.
\end{proof}

\subsection{An explicit simple-root formula}

Suppose $f_0(a)=0$, $f_0'(a)\ne0$, and
$u(z)=f_0(z)/(z-a)$ is nonzero on a small disc. For
$F=f_0+R$ with $R$ positive-supported, the unique zero $a+\eta$ in the monad
has the exact expansion
\begin{equation}\label{eq:lagrange}
 \eta=\sum_{n\ge1}\frac{(-1)^n}{n!}
 \left.\frac{\mathrm d^{n-1}}{\mathrm dz^{n-1}}
 \left(\frac{R(z)}{u(z)}\right)^n\right|_{z=a}.
\end{equation}

\begin{proof}
Put $\phi(X)=-R(a+X)/u(a+X)$, introduce a formal parameter $\lambda$,
and solve $X=\lambda\phi(X)$ recursively in $\K[[\lambda]]$.
For clarity, the required formal Lagrange identity follows directly from
formal residues. If $\phi(0)\ne0$, the substitution
$\lambda=X/\phi(X)$ is an invertible formal change of variable. Therefore,
for $n\ge1$,
\begin{align*}
 [\lambda^n]X(\lambda)
 &=\Res_{X=0}\left(
 \frac{\phi(X)^n}{X^n}
 -\frac{\phi(X)^{n-1}\phi'(X)}{X^{n-1}}\right)\dd X\\
 &=\frac1n[X^{n-1}]\phi(X)^n.
\end{align*}
In the second equality, use $(\phi^n)'=n\phi^{n-1}\phi'$ and compare
coefficients. Invariance of the formal residue under an invertible change
of variable follows by checking Laurent monomials: derivatives have zero
residue, and the transformed logarithmic differential has residue one.
If $\phi(0)=0$, the unique solution is $X=0$, and the same coefficient
identity is zero on both sides because $\phi^n$ is divisible by $X^n$.

Thus the coefficient of $\lambda^n$ is the $n$th summand in
\eqref{eq:lagrange}. Every such summand contains $n$ factors from the
positive support of $R$. Neumann's lemma
therefore permits the substitution $\lambda=1$ as a Hahn sum. The formal
equation becomes $f_0(a+\eta)+R(a+\eta)=0$. Uniqueness follows either from
\cref{thm:zerolifting} with $m=1$ or from the difference factorization above.
\end{proof}

\subsection{Several variables}

The same definitions work with $U\subseteq\C^d$, multi-index Taylor series,
and a common well-ordered support. Suppose an ordinary holomorphic equation
$H_0(x,y)=0$ has a holomorphic solution $y=g_0(x)$ on an ordinary domain and
its $y$-Jacobian is invertible along that solution. If
$H=H_0+R$ is a positive-supported Hahn perturbation, there is a unique
positive-supported correction $g=g_0+S$ solving $H(x,g(x))=0$.
At each exponent the recursion is
\[
 \frac{\partial H_0}{\partial y}(x,g_0(x))S_\gamma(x)
   =-\text{the already determined coefficient at exponent }\gamma.
\]
The matrix inverse is holomorphic, and the support lemma makes every
right-hand side finite. This proves the multi-variable statement, including
actual evaluation at infinitesimal tuples. It is a concrete finite-dimensional
instance of the general support-based implicit-function philosophy in
\cite{vdH}.

\subsection{Local mapping at arbitrary surcomplex points}

The leading coefficient need not have a simple zero for a useful local
mapping theorem. Rescaling an arbitrarily small neighborhood recovers a
polynomial leading term, even at a point with infinitesimal derivative.

\begin{theorem}[Local finite-degree mapping and open mapping]\label{thm:openmapping}
Let $U$ be connected, let $F\in\HH(U)$ be nonconstant, and fix
$\zeta\in U^\sharp$. Write
\[
 A_n=\frac{F^{(n)}(\zeta)}{n!},
 \qquad m=\min\{n\ge1:A_n\ne0\}.
\]
There are arbitrarily small positive monomials $\rho$ such that
\begin{equation}\label{eq:localnormalization}
 G(X)=\frac{F(\zeta+\rho X)-F(\zeta)}{A_m\rho^m}
      =X^m+R(X),\qquad R\in\HH(\C)
\end{equation}
has positive support. For each such $\rho$,
\[
 F(\zeta+\rho\OO)=F(\zeta)+A_m\rho^m\OO,
\]
and every fiber in these neighborhoods has exactly $m$ points counted with
analytic multiplicity. Consequently $F$ is an open map for the fine
topology. If $F'(\zeta)\ne0$, it has an analytic local inverse in affine
Hahn-supported charts, with derivative $1/F'$ at corresponding points.
\end{theorem}
\begin{proof}
The integer $m$ exists by \cref{thm:identity}, applied to $F-F(\zeta)$.
Write $\zeta=a+\epsilon$. The supports of all $A_n$ lie in the same
well-ordered set
\[
 T=S+M(\supp\epsilon),
\]
where $S$ supports $F$, and $\val(A_n)\ge s=\min S$.
After division by $A_m$, their supports lie in the well-ordered set
$T+\supp(A_m^{-1})$, bounded below by $s-\val(A_m)$.
Choose a positive monomial $\rho=t^\delta$ with
\[
 \delta>0,\qquad \delta>\val(A_m)-s.
\]
Taylor's theorem gives the candidate
\[
 G(X)=X^m+\sum_{n>m}(A_n/A_m)\rho^{n-m}X^n.
\]
Neumann's lemma makes this jointly summable: the coefficient supports are
contained in $T+\supp(A_m^{-1})+\{k\delta:k\ge1\}$, and only finitely
many $n$ contribute at a fixed exponent. Thus each resulting coefficient
function is a polynomial in $X$. Every nonleading exponent is positive by
the choice of $\delta$, proving \eqref{eq:localnormalization} as an element
of $\HH(\C)$. The equality with the actual function follows from Taylor
evaluation and the same joint support calculation, for every finite $X$.

For $b\in\OO$, the reduction of $G-b$ is the ordinary polynomial
$X^m-\st b$. Its ordinary zeros have total multiplicity $m$. All zeros of
$G-b$ in $\OO$ must reduce to one of them; \cref{thm:zerolifting} gives
exactly the corresponding number of roots in each monad. This proves the
image and fiber assertions. Affine rescaling preserves multiplicities.

Given any fine neighborhood of $\zeta$, increase $\delta$ further so that
$\zeta+\rho\OO$ is contained in that neighborhood. This is possible by
requiring $\delta$ to exceed the valuation of a ball radius. Its image is
$F(\zeta)+A_m\rho^m\OO$, a fine-open neighborhood. Applying this at
every point proves the open-mapping assertion. If $m=1$, apply
\cref{thm:inverse} to $G=X+R$ on $\C$ and then undo the two affine
rescalings. The chain rule gives the derivative formula.
\end{proof}

The finite-degree assertion is a multiplicity statement about fibers, not a
claim that these totally disconnected neighborhoods form a classical
branched covering of surfaces. It recovers the local mapping consequences
of complex analyticity without importing an inappropriate surface topology.

\section{Contours, Cauchy's formula, and Morera's theorem}\label{sec:contours}

\subsection{Definition of the contour integral}

Let $\sigma$ be an ordinary piecewise $C^1$ complex curve, and let
$H=\sum h_\gamma t^\gamma$ have one well-ordered support with continuous
complex coefficients along the curve. Define
\begin{equation}\label{eq:contourdefinition}
 \cint_\sigma H(w)\dd w
 :=\sum_\gamma t^\gamma\int_\sigma h_\gamma(w)\dd w.
\end{equation}
The output is a surcomplex number because its support is contained in the
original support. Equality of coefficientwise expressions on the curve is
pointwise equality of normal forms, so the integral does not depend on adding
zero coefficients or on a redundant choice of support.

This is a \emph{definition}, not an exchange of a limit with an infinite sum.
It is $\K$-linear, since multiplying by a fixed Hahn series produces a
summable convolution with finitely many contributions at each exponent.
It has ordinary orientation, subdivision, and reparametrization properties.
It is not the integral of a fine-continuous map from a real interval; such a
map would be constant by \cref{cor:paths}.

\begin{theorem}[Cauchy theorem and primitives]\label{thm:cauchytheorem}
If $U$ is simply connected and $F\in\HH(U)$, then every closed ordinary
piecewise $C^1$ curve $\sigma\subset U$ satisfies
$\cooint_\sigma F(w)\dd w=0$.
There exists $A\in\HH(U)$ with $A'=F$. Given $a\in U$, a unique such
primitive satisfies $A(a)=0$, and
\[
 A(b)-A(a)=\cint_\sigma F(w)\dd w
\]
for every ordinary curve from $a$ to $b$. Fine differentiation interprets
$A'=F$ at every point of $U^\sharp$.
\end{theorem}
\begin{proof}
Choose the ordinary holomorphic primitive of each $f_\gamma$, normalized to
vanish at $a$. The resulting series has the same support and differentiates
to $F$. Cauchy's theorem and the endpoint formula hold at each coefficient.
Uniqueness follows because a coefficientwise derivative-zero function on a
connected domain has constant coefficients.
\end{proof}

\subsection{Cauchy's integral formula at nonstandard points}

\begin{theorem}[Surcomplex Cauchy formula]\label{thm:cauchyformula}
Let $D$ be a bounded Jordan domain with piecewise $C^1$ boundary $\sigma$,
and let $\overline D\subset U$. If $F\in\HH(U)$, $z\in D^\sharp$, and
$k\ge0$, then
\begin{equation}\label{eq:cauchyformula}
 \frac{F^{(k)}(z)}{k!}
 =\frac1{2\pi i}\cooint_\sigma
      \frac{F(w)}{(w-z)^{k+1}}\dd w.
\end{equation}
The integral on the right is well defined by coefficientwise expansion of
the kernel.
\end{theorem}
\begin{proof}
Write $z=a+\epsilon$ with $a\in D$. On the ordinary curve,
\begin{equation}\label{eq:kernel}
 (w-z)^{-k-1}
 =\sum_{n\ge0}\binom{n+k}{k}
    \frac{\epsilon^n}{(w-a)^{n+k+1}}.
\end{equation}
Because $w-a$ is a nonzero ordinary complex number and is bounded away from
zero on the curve, these coefficients are continuous. The fully expanded
support of the product with $F(w)$ lies in
$S+M(\supp\epsilon)$, with finite multiplicity at each exponent. Hence the
integral exists and coefficientwise integration can be grouped as written.
Ordinary Cauchy's formula gives
\[
 \frac1{2\pi i}\int_\sigma
 \frac{f_\gamma(w)}{(w-a)^{n+k+1}}\dd w
 =\frac{f_\gamma^{(n+k)}(a)}{(n+k)!}.
\]
Multiplying by the binomial coefficient leaves
$f_\gamma^{(n+k)}(a)/(k!n!)$. Summing recovers
$F^{(k)}(a+\epsilon)/k!$ by Taylor realization.
\end{proof}

\subsection{Coefficientwise estimates}

For a standard closed circle $|w-a|=r$ contained with its interior in $U$, define
\[
 M_\gamma(r)=\max_{|w-a|=r}|f_\gamma(w)|,
 \qquad M(r)=\sum_\gamma M_\gamma(r)t^\gamma.
\]
The latter is a nonnegative surreal number. The usual estimates imply
\begin{equation}\label{eq:cauchyestimates}
 |F^{(n)}(a)|\le \frac{n!}{r^n}M(r).
\end{equation}
For completeness, the absolute-coefficient majorant is legitimate: if
$Z=\sum c_\gamma t^\gamma$ and $B=\sum |c_\gamma|t^\gamma$, then
$B^2-Z\overline Z$ has nonnegative real coefficients, since each real
contribution is bounded by the corresponding product of absolute values.
Thus $|Z|\le B$. Apply this to the ordinary coefficientwise Cauchy estimates.

One must not let a standard radius tend to infinity in
\eqref{eq:cauchyestimates} as though it were cofinal among all surreal
radii. That error would produce a false Liouville theorem.

\begin{theorem}[Morera]\label{thm:morera}
Let $H=\sum h_\gamma t^\gamma$ have one well-ordered support and continuous
complex coefficients on an ordinary open set $U$. If
\[
 \cooint_{\partial T}H(w)\dd w=0
\]
for every closed triangle $T$ with $T\subset U$, then all $h_\gamma$ are
holomorphic. Hence $H$ has the canonical realization in $\HH(U)$.
\end{theorem}
\begin{proof}
Uniqueness of normal forms makes each ordinary coefficient integral zero.
Ordinary Morera's theorem applies to each continuous coefficient separately.
\end{proof}

\section{Meromorphic functions and actual surcomplex residues}\label{sec:residues}

\subsection{Quotients and local poles}

A meromorphic object in this framework is a quotient $M=F/G$ with
$F,G\in\HH(U)$ and $G\ne0$. On any ordinary subdomain on which the leading
coefficient of $G$ is nonzero, \cref{prop:units} expands the quotient into
$\HH$ there. Near a zero of that leading coefficient, preparation gives the
proper local description.

\begin{proposition}[Analytic part plus a finite rational part]\label{prop:localrational}
Let $G=t^s(g_0+R)$ and suppose $g_0$ has a zero of order $m$ at $a$.
On a sufficiently small ordinary disc $D$ about $a$, there are
$Q\in\HH(D)$ and polynomials $A,P\in\K[X]$, with $P$ monic of degree $m$
and $\deg A<m$, such that
\begin{equation}\label{eq:localrational}
 \frac{F(z)}{G(z)}=Q(z)+\frac{A(z-a)}{P(z-a)}.
\end{equation}
All roots of $P$ are infinitesimal. In particular, there are only finitely
many poles in $a+\mm$, after cancellations.
\end{proposition}
\begin{proof}
Prepare $G=t^s u_0 V P$. The factor $t^s u_0V$ is invertible in $\HH(D)$.
Set $H=F/(t^s u_0V)$ and divide $H$ by $P$ using
\cref{thm:division}. This gives \eqref{eq:localrational}. The assertion about
roots is \cref{thm:zerolifting}.
\end{proof}

Over the algebraically closed field $\K$, partial fractions write the rational
part uniquely as
\[
 \frac{A(X)}{P(X)}=
 \sum_{\eta}\sum_{j=1}^{m_\eta}\frac{c_{\eta,j}}{(X-\eta)^j}.
\]
Define the actual residue at the pole $\zeta=a+\eta$ by
\begin{equation}\label{eq:actualresidue}
 \Res_{z=\zeta}M(z)\dd z=c_{\eta,1}.
\end{equation}
This is independent of the local decomposition: after subtracting the finite
principal parts, the remainder is analytic at $\zeta$; uniqueness of local
Taylor coefficients gives uniqueness of the coefficient of
$(z-\zeta)^{-1}$. It is also the residue of the local Laurent germ in a
sufficiently small punctured fine neighborhood of $\zeta$.

\begin{theorem}[Removable singularities in the meromorphic category]\label{thm:removable}
A meromorphic quotient as above has a removable singularity at an actual
surcomplex point $\zeta$ if and only if it is bounded by some fixed positive
surreal on a sufficiently small punctured fine neighborhood of $\zeta$.
\end{theorem}
\begin{proof}
The local rational decomposition gives a finite Laurent principal part.
If there are no negative powers, the local Taylor representation extends
across $\zeta$; continuity bounds its modulus on a sufficiently small ball.
If a pole of order $k>0$ remains, then
\[
 M(\zeta+h)=h^{-k}(c+o(1)),\qquad c\ne0.
\]
The expression in parentheses has modulus greater than $|c|/2$ for all
sufficiently small nonzero $h$. Choosing $|h|$ still smaller makes the
modulus exceed any prescribed positive surreal bound. A pole is therefore
incompatible with local boundedness.
\end{proof}

This does not extend to arbitrary fine-holomorphic punctured germs. For
example, on $\OO\setminus\{0\}$ let $B(h)=1$ when $\val(h)$ is an ordinal
and $B(h)=0$ otherwise. Valuation is locally constant at nonzero points, so
$B$ is locally constant and fine-holomorphic, and $|B|\le1$. However, for
arbitrarily large ordinals $\beta$,
\[
 B(t^\beta)=1,\qquad B(t^{\beta+1/2})=0.
\]
Both arguments can be placed in any prescribed punctured ball about zero.
Thus $B$ has no continuous extension there. The coherent meromorphic
hypothesis in \cref{thm:removable} cannot simply be deleted.

\subsection{The contour theorem}

\begin{theorem}[Residue theorem at surcomplex poles]\label{thm:residues}
Let $D\Subset U$ be a bounded Jordan domain with piecewise $C^1$ boundary
$\sigma$. Let $M=F/G$, where $F,G\in\HH(U)$ and the leading coefficient
$g_0$ of the normalized denominator is nonzero on $\sigma$. Then $M$ has
finitely many poles in $D^\sharp$, and
\begin{equation}\label{eq:residuetheorem}
 \cooint_\sigma M(z)\dd z
 =2\pi i\sum_{\zeta\in\operatorname{Pole}(M)\cap D^\sharp}
       \Res_{z=\zeta}M(z)\dd z,
\end{equation}
where the displayed sum is over the finite set of actual poles.
\end{theorem}
\begin{proof}
There are finitely many zeros $a_1,\ldots,a_r$ of $g_0$ in $D$, since
$g_0$ is holomorphic near $\overline D$ and nonzero on its boundary.
Outside their monads, reduction shows that $G$ is nonzero. Near each $a_j$,
\cref{prop:localrational} gives finitely many poles. This proves finiteness.

On $U\setminus\{g_0=0\}$, expand $G^{-1}$ by its positive-support geometric
series. Its coefficients are ordinary meromorphic functions with possible
poles only at the zeros of $g_0$; at every fixed exponent only finitely many
powers occur. The ordinary residue theorem therefore reduces the outer
contour, coefficient by coefficient, to a sum of small ordinary circles about
the $a_j$.

On one such circle use \eqref{eq:localrational}. The $Q$ integral is zero.
For an infinitesimal $\eta$, the expansion
\[
 (X-\eta)^{-j}=\sum_{n\ge0}\binom{j+n-1}{n}\eta^nX^{-j-n}
\]
is jointly summable on the ordinary circle. Its coefficientwise integral is
$2\pi i$ for $j=1$ and $0$ for $j>1$, because an $X^{-1}$ term occurs only
when $j=1,n=0$. Thus the small-circle integral equals $2\pi i$ times the
sum of the actual residues at the poles in $a_j+\mm$. Summing proves
\eqref{eq:residuetheorem}.
\end{proof}

\begin{corollary}[Argument principle]\label{cor:argument}
Suppose the leading coefficients of nonzero $F,G\in\HH(U)$ are nonzero
on $\sigma$, and let $M=F/G$. Then
\begin{equation}\label{eq:argument}
 \frac1{2\pi i}\cooint_\sigma\frac{M'(z)}{M(z)}\dd z
 =\sum_{\zeta\in Z\cup P}\ord_\zeta M
 =N_F(D^\sharp)-N_G(D^\sharp).
\end{equation}
Here $Z$ and $P$ are the finite sets of actual zeros and poles of $M$ in
$D^\sharp$. Multiplicities are finite integers and cancellations are
respected.
\end{corollary}
\begin{proof}
At a point of order $m\in\Z$, write locally
$M(z)=(z-\zeta)^mU(z)$ with $U(\zeta)\ne0$. Then
$M'/M=m/(z-\zeta)+U'/U$ has residue $m$. Apply
\cref{thm:residues} to $M'/M$, a quotient of Hahn-supported holomorphic
functions. The local preparations ensure there are only finitely many zeros
and poles in the indicated halo.
\end{proof}

The left side is therefore an ordinary integer even when individual zeros,
poles, or residues involve infinitesimal and infinite surreal quantities.
It is not a winding number of a fine-continuous real loop; its meaning is the
specified contour pairing and the divisor identity.

\subsection{A base-independent rational residue identity}

For a rational function $R(z)\in\K(z)$, residues make sense at all poles
in the affine surcomplex plane, including infinitely large points, without
choosing a shadow domain. Define the residue at
infinity using $w=1/z$:
\[
 \Res_\infty R(z)\dd z=\Res_{w=0}\bigl(-w^{-2}R(1/w)\bigr)\dd w.
\]
Partial fractions immediately prove
\begin{equation}\label{eq:rationalglobal}
 \sum_{a\in\operatorname{Pole}(R)}\Res_{z=a}R(z)\dd z
 +\Res_\infty R(z)\dd z=0,
\end{equation}
where $\operatorname{Pole}(R)\subset\K$ is the finite set of affine
poles. Indeed, a simple-pole term contributes its coefficient at its
pole and its negative at infinity; higher-pole and polynomial terms contribute
zero total. A rational differential has a rational primitive if and only if
all its affine residues vanish, because every remaining partial fraction
integrates rationally in characteristic zero.

\section{Rouch\'e's theorem and exact conservation of zeros}\label{sec:rouche}

\begin{theorem}[Reduction controls all zeros]\label{thm:rouche}
Let $F=t^s(f_0+R)\in\HH(U)$, with $R$ positive-supported and $f_0\ne0$.
Let $D\Subset U$ be a bounded Jordan domain whose boundary contains no zero
of $f_0$. Then
\begin{equation}\label{eq:zerocount}
 N_F(D^\sharp)=N_{f_0}(D).
\end{equation}
For each zero $a$ of $f_0$, exactly $\ord_a f_0$ of these zeros lie in
$a+\mm$, counted with multiplicity.
\end{theorem}
\begin{proof}
If $F(z)=0$ and $z\in D^\sharp$, reduction of $t^{-s}F(z)$ gives
$f_0(\st z)=0$. Thus every zero lies in the monad of an ordinary zero of
$f_0$. There are finitely many such zeros in $D$. Apply
\cref{thm:zerolifting} in disjoint small discs about them and sum the
multiplicities.
\end{proof}

\begin{corollary}[Infinitesimal Rouch\'e theorem]\label{cor:infrouche}
If $F=t^s(f_0+R)$ and $G=t^s(f_0+S)$ have the same leading coefficient and
$R,S$ are positive-supported, then they have the same number of zeros in
$D^\sharp$ whenever $f_0$ is nonzero on $\partial D$. This remains true
regardless of how their multiple zeros split within individual monads.
\end{corollary}

There is also a version closer to the usual inequality.

\begin{corollary}[Rouch\'e with a real margin]\label{cor:marginrouche}
Suppose $t^{-s}F$ and $t^{-s}G$ are nonnegative-supported, with reductions
$f_0$ and $g_0$, and $f_0$ is nonzero on $\partial D$. If, for one ordinary
real $0\le q<1$,
\begin{equation}\label{eq:rouchemargin}
 |G(a)-F(a)|\le q|F(a)|\qquad(a\in\partial D),
\end{equation}
then $N_F(D^\sharp)=N_G(D^\sharp)$.
\end{corollary}
\begin{proof}
Divide the inequality by the positive scalar $t^s$ and take standard parts.
This gives $|g_0-f_0|\le q|f_0|<|f_0|$ on the boundary. Ordinary
Rouch\'e gives equal counts and ensures $g_0$ is nonzero there. Now use
\cref{thm:rouche} for both functions.
\end{proof}

A strict surreal inequality with no real margin cannot simply be reduced to
a strict complex inequality: strictness can disappear under standard part.
The infinitesimal version and the real-margin version state exactly the
conditions used in their proofs.

\section{Maximum modulus, Liouville, and Schwarz: the correct distinctions}\label{sec:growth}

\subsection{The local maximum principle survives}

\begin{theorem}[Maximum-modulus principle]\label{thm:maximum}
Let $U$ be connected and $F\in\HH(U)$. If $|F|$ has a local maximum at a
point of $U^\sharp$ for the fine topology, then $F$ is constant on
$U^\sharp$.
\end{theorem}
\begin{proof}
Suppose $F$ is not constant and consider a point $z$. By
\cref{thm:identity} applied to $F-F(z)$, there is a least $m\ge1$ for which
$A_m=F^{(m)}(z)/m!\ne0$. Write $c=F(z)$. If $c=0$, isolation of zeros
already rules out a local maximum.

If $c\ne0$, choose $u\in\K$ with $|u|=1$ such that
$A_mu^m/c=|A_m|/|c|>0$. Such a phase exists by algebraic closure: take an
$m$th root of the unit $c|A_m|/(|c|A_m)$.
For arbitrarily small positive surreal $r$, put $h=ru$. The Taylor formula
gives
\[
 \frac{F(z+h)}c=1+\delta+E,
 \qquad \delta=\frac{|A_m|}{|c|}r^m>0.
\]
The remainder estimate allows $r$ to be chosen so that $|E|\prec\delta$.
Consequently $\operatorname{Re}(F(z+h)/c)>1$, hence $|F(z+h)|>|c|$.
Such $h$ can be made smaller than any proposed neighborhood radius, a
contradiction.
\end{proof}

\subsection{Boundedness is not coefficientwise boundedness}

\begin{example}[A bounded nonconstant function on the entire finite halo]\label{ex:liouville}
For $t=\omega^{-1}$, the function $F(z)=tz$ belongs to $\HH(\C)$ and is
nonconstant. Yet $|F(z)|<1$ for every $z\in\C^\sharp=\OO$, because a
finite element times a positive infinitesimal is infinitesimal. Thus even a
fixed real bound on the whole finite halo does not imply constancy.
The function $z$ itself is bounded there by the infinite bound $\omega$.
\end{example}

\begin{theorem}[Coefficientwise Liouville theorem]\label{thm:liouville}
If $F=\sum f_\gamma t^\gamma\in\HH(\C)$ and every entire coefficient
$f_\gamma$ is bounded on $\C$ by some ordinary real constant depending on
$\gamma$, then $F$ is constant. More generally, if all $f_\gamma$ have
polynomial growth of degree at most one fixed integer $d$, then $F$ is a
polynomial in $z$ of degree at most $d$ with coefficients in $\K$.
\end{theorem}
\begin{proof}
Ordinary Cauchy estimates with real radii tending to infinity show that each
$f_\gamma$ is constant in the bounded case. Under degree-$d$ growth, those
estimates show $f_\gamma^{(d+1)}=0$. Collect the finitely many powers
$1,z,\ldots,z^d$ across the common Hahn support to obtain a polynomial over
$\K$.
\end{proof}

\subsection{Schwarz's exact estimate need not survive}

Let $\mathbb D$ denote the ordinary unit disc. Its halo is
$\mathbb D^\sharp=\{z\in\OO:|\st z|<1\}$, not the entire surreal-radius
ball $\{z:|z|<1\}$. The latter additionally contains some points with
standard-part modulus exactly $1$.

\begin{example}[Failure of the exact Schwarz inequality]\label{ex:schwarz}
Set $F(z)=(1+t)z$ on $\mathbb D^\sharp$. It is in $\HH(\mathbb D)$,
$F(0)=0$, and $|F(z)|<1$ throughout the halo: each standard part has a
positive ordinary margin from the unit circle. Nevertheless
\[
 F'(0)=1+t>1,\qquad |F(z)|=(1+t)|z|>|z|\quad(z\ne0).
\]
The classical Schwarz estimate is therefore false for these halo domains.
\end{example}

\begin{proposition}[Shadow Schwarz theorem]\label{prop:schwarz}
Let $F\in\HH(\mathbb D)$ be nonnegative-supported, suppose $F(0)=0$, and
assume $|F(z)|\le1$ for every $z\in\mathbb D^\sharp$. Its reduction $f_0$
satisfies the ordinary Schwarz lemma. In particular,
\[
 |\st F'(0)|\le1,
 \qquad \st\left(\left|\frac{F(z)}z\right|\right)\le1
 \quad(0\ne z\in\mathbb D^\sharp).
\]
The quotient in the second formula is always finite.
\end{proposition}
\begin{proof}
For standard arguments, reduction gives $|f_0(a)|\le1$ and $f_0(0)=0$.
Apply the ordinary Schwarz lemma. If $\st z=a\ne0$, reduction of $F(z)/z$
is $f_0(a)/a$. If $z$ is infinitesimal, $F(0)=0$ and Taylor factorization
give $F(z)/z=F'(0)+zB(z)$ with nonnegative value; its standard part is
$f_0'(0)$. This proves both finiteness and the bound.
\end{proof}

The examples do not contradict the maximum principle: these functions never
attain a local maximum. Nor do they refute versions on other non-Archimedean
analytic domains equipped with different boundary or boundedness conditions.

\section{Arbitrary centers and scales}\label{sec:charts}

The finite halo is a natural starting domain, not a restriction to finite
surcomplex locations. Let $A\in\K$ and $B\in\K^\times$. For an ordinary
complex domain $U$, put
\[
 \Omega_{A,B}(U)=A+B U^\sharp.
\]
Given $F\in\HH(U)$, define
$\mathcal F(z)=F((z-A)/B)$ on this affine chart. All local conclusions
transport exactly, with
\begin{equation}\label{eq:affinederivative}
 \mathcal F^{(n)}(z)=B^{-n}F^{(n)}((z-A)/B).
\end{equation}
Thus the center $A$ can have infinite real and imaginary parts, and $B$ may
be arbitrarily large or small.

For the parametrized shadow contour $A+B\sigma$, define
\begin{equation}\label{eq:affineintegral}
 \cint_{A+B\sigma}\mathcal F(z)\dd z
 :=B\cint_\sigma F(w)\dd w.
\end{equation}
This is an affine change-of-variable convention. Residues of differentials
are invariant: if $z-\zeta=B(w-\eta)$, the factor $B$ in $\dd z$ cancels
the factor $B^{-1}$ in a simple pole. Cauchy, preparation, zero counting, and
the residue theorem therefore hold in every such chart.

Every surcomplex point lies in a chart of this form. This does not force
global continuation between arbitrary charts. An atlas requires compatible
expressions on overlaps; merely knowing the local derivative is insufficient.
In particular, one must not apply the connectedness of the ordinary plane to
the fine topology of all surcomplex points.

\section{Exponentials, logarithms, and unavoidable nonstandard periods}\label{sec:globalexp}

This section gives a genuinely global construction, while also displaying a
sharp obstruction to copying the classical exponential covering theory.
We use the established surreal exponential $\exp_{\No}$ and its inverse
$\log_{\No}$: it is an ordered-group isomorphism from $(\No,+)$ to
$(\No_{>0},\cdot)$, extends the real exponential, and agrees with its Taylor
series on real infinitesimals \cite{Gonshor,BMderivations}.

\subsection{The infinitesimal exponential and logarithm}

For $h\in\mm$ and $u\in\mm$, define
\begin{equation}\label{eq:small-exp-log}
 e_{\mm}(h)=\sum_{n\ge0}\frac{h^n}{n!},\qquad
 \ell_{\mm}(1+u)=\sum_{n\ge1}\frac{(-1)^{n+1}}n u^n.
\end{equation}
Formal identities and \cref{prop:formal} give mutually inverse group
isomorphisms
\[
 e_{\mm}:(\mm,+)\longrightarrow(1+\mm,\cdot),
 \qquad \ell_{\mm}=e_{\mm}^{-1}.
\]
For example, the exponential addition law follows by the binomial theorem in
a jointly summable double series. The logarithmic inverse identities are
formal identities with zero substituted constant term. These operations
commute with conjugation.

\subsection{A canonical exponential on the finite-imaginary strip}

Let $\OO_\R=\OO\cap\No$ and put
\[
 \mathcal S=\No+i\OO_\R.
\]
For finite real $y=b+\epsilon$ with $b\in\R$ and
$\epsilon\in\mm\cap\No$, set
\[
 E_{\mathrm{ang}}(y)=e^{ib}e_{\mm}(i\epsilon).
\]
Define
\begin{equation}\label{eq:stripexp}
 \Exp_{\mathcal S}(x+iy)=\exp_{\No}(x)E_{\mathrm{ang}}(y).
\end{equation}

\begin{theorem}[Surjective strip exponential]\label{thm:stripexp}
The map $\Exp_{\mathcal S}$ is a surjective group homomorphism
$\mathcal S\to\K^\times$. It extends the ordinary complex exponential and
the real surreal exponential, commutes with conjugation, and satisfies
\begin{equation}\label{eq:stripexp-properties}
 |\Exp_{\mathcal S}(x+iy)|=\exp_{\No}(x),\qquad
 \ker\Exp_{\mathcal S}=2\pi i\Z.
\end{equation}
For every $z\in\mathcal S$ and infinitesimal $h\in\K$,
\[
 \Exp_{\mathcal S}(z+h)=\Exp_{\mathcal S}(z)e_{\mm}(h).
\]
Consequently it is locally strongly analytic and its fine derivative equals
itself.
\end{theorem}
\begin{proof}
The angular addition law follows by splitting finite real arguments into
standard and infinitesimal parts. The real exponential law gives the full
group law. Conjugation commutes with both factors. The angular factor has
modulus $1$, because its product with its conjugate is $1$.

If the exponential is $1$, its modulus forces $x=0$. Reduction gives
$e^{ib}=1$, so $b\in2\pi\Z$. Then $e_{\mm}(i\epsilon)=1$, and the
infinitesimal logarithm forces $\epsilon=0$. This proves the kernel claim.

For surjectivity, let $w\ne0$, set $r=|w|$, and $u=w/r$. Then $|u|=1$,
so $u$ is finite and $u_0=\st u$ lies on the ordinary unit circle. Choose a
real angle $b$ with $e^{ib}=u_0$. The element $q=e^{-ib}u$ belongs to
$1+\mm$ and satisfies $\overline q=q^{-1}$. Hence
\[
 \overline{\ell_{\mm}(q)}
 =\ell_{\mm}(\overline q)=\ell_{\mm}(q^{-1})=-\ell_{\mm}(q).
\]
Thus $\ell_{\mm}(q)=i\epsilon$ for an infinitesimal real $\epsilon$.
It follows that
$w=\Exp_{\mathcal S}(\log_{\No}r+i(b+\epsilon))$.

Finally, for $h=p+iq\in\mm$, the exponential laws and the real
infinitesimal Taylor expansion give a local multiplier
$e_{\mm}(p)e_{\mm}(iq)=e_{\mm}(h)$. Its first-order Taylor estimate gives
the derivative, and the complete expansion gives strong local analyticity.
\end{proof}

Thus every nonzero surcomplex number has logarithms with finite imaginary
part, uniquely modulo $2\pi i\Z$. This assertion needs no choice of a sine
function at an infinite real argument.

\subsection{Explicit exponential extensions to all of the plane}

The normal form of a real surreal $y$ splits uniquely as
\[
 y=y_\infty+y_{\fin},
\]
where $y_\infty$ contains the monomials $\omega^a$ with $a>0$, and
$y_{\fin}$ contains the real and infinitesimal terms. Both summands are
surreal numbers and the projection $y\mapsto y_{\fin}$ is additive.
Let $c_\omega(y)\in\R$ be the coefficient of the monomial $\omega$ in $y$.
For a real parameter $\lambda$, put
\begin{equation}\label{eq:phaseprojection}
 p_\lambda(y)=y_{\fin}+\lambda c_\omega(y)\in\OO_\R,
 \qquad
 \Exp_\lambda(x+iy)=\exp_{\No}(x)E_{\mathrm{ang}}(p_\lambda(y)).
\end{equation}

\begin{theorem}[Global entire exponential extensions]\label{thm:globalexp}
Each $\Exp_\lambda:\K\to\K^\times$ is a surjective, conjugation-compatible
group homomorphism extending $\Exp_{\mathcal S}$. It is locally represented
at every point by the exact series
\begin{equation}\label{eq:globalexp-taylor}
 \Exp_\lambda(z+h)=\Exp_\lambda(z)
          \sum_{n\ge0}\frac{h^n}{n!}\qquad(h\in\mm),
\end{equation}
and therefore $\Exp_\lambda'=\Exp_\lambda$ in the fine sense. However,
\[
 \Exp_\lambda(i\omega)=e^{i\lambda}.
\]
The global exponential is consequently not determined by its values on the
finite-imaginary strip, by conjugation, or by its local differential equation.
For $\lambda=0$,
\begin{equation}\label{eq:global-kernel}
 \ker\Exp_0=i\bigl(\No_\infty+2\pi\Z\bigr),
\end{equation}
where $\No_\infty$ is the additive subgroup of purely infinite real normal
forms, including zero.
\end{theorem}
\begin{proof}
The map $p_\lambda$ is additive and fixes every finite real surreal. The
angular and real exponential laws prove the group law and extension
properties. Surjectivity follows already from the strip. Conjugation follows
from additivity under negation.

For an infinitesimal real $q$, $c_\omega(q)=0$ and $p_\lambda(q)=q$.
Hence the change in the angular projection under an infinitesimal increment
is the increment itself, giving \eqref{eq:globalexp-taylor} exactly as in the
strip case. The value at $i\omega$ is immediate. For $\lambda=0$, the
modulus again forces the real part of a kernel element to be zero, and the
strip kernel condition says precisely $y_{\fin}\in2\pi\Z$.
\end{proof}

These formulas exhibit a particular normal-form phase convention; they are
not claimed to be the unique or preferred global surcomplex exponential.
They also do not identify their enormous kernel with the ring of omnific
integers. Different axioms for global phases can select other functions.
Although two such exponentials agree on a fine-open strip, they can disagree
elsewhere. That fact is compatible with \cref{thm:identity}, which requires
one coherent coefficient domain rather than arbitrary global local germs.

\subsection{A no-go theorem for the classical period group}

\begin{theorem}[Every global extension has nonstandard periods]\label{thm:periodobstruction}
Suppose $E:(\K,+)\to(\K^\times,\cdot)$ is a group homomorphism which
extends $\Exp_{\mathcal S}$ and satisfies
$E(\overline z)=\overline{E(z)}$. Then, for every infinite real surreal $y$,
there is a finite real surreal $v$ such that
\begin{equation}\label{eq:periodobstruction}
 i(y-v)\in\ker E.
\end{equation}
In particular, no such global extension can have kernel exactly
$2\pi i\Z$.
\end{theorem}
\begin{proof}
The homomorphism and conjugation identities imply
$E(iy)\overline{E(iy)}=E(iy)E(-iy)=1$. Thus $E(iy)$ lies on the surcomplex
unit circle. The angular part of \cref{thm:stripexp} parametrizes the whole
unit circle using finite real angles. Choose finite $v$ with
$E(iy)=\Exp_{\mathcal S}(iv)=E(iv)$. Then $E(i(y-v))=1$.
Since $y-v$ is infinite, this period is not an ordinary integer multiple of
$2\pi$.
\end{proof}

The obstruction is algebraic, not merely topological. Extending the ordinary
complex exponential to the whole surcomplex plane while keeping conjugation
and the finite-angle values necessarily changes its period structure.

\subsection{Local logarithms and local conformality}

Choose $w\ne0$ and a strip logarithm $L$ with
$\Exp_{\mathcal S}(L)=w$. On $w(1+\mm)$ define
\begin{equation}\label{eq:locallog}
 \Log_w(w+h)=L+\ell_{\mm}(1+h/w).
\end{equation}
For any extension $\Exp_\lambda$, this is an inverse of its restriction to
$L+\mm$, and
\[
 \Log_w'(z)=1/z.
\]
Both assertions follow from the inverse formal series and its derivative.
Changing $L$ by $2\pi i n$ changes the logarithm by that constant. Thus
familiar local logarithmic calculus survives even though a classical global
covering-space interpretation does not.

\section{Worked examples and reproducible checks}\label{sec:examples}

\subsection{A simple zero moved by an exponential perturbation}

Consider
\[
 F(z)=z+t e^z\in\HH(\C),\qquad t=\omega^{-1}.
\]
Its leading coefficient is $z$. It has exactly one zero in $\mm$, and no
other zero anywhere in the finite halo. Formula \eqref{eq:lagrange} gives
\begin{equation}\label{eq:Wexample}
 z_* =\sum_{n\ge1}\frac{(-1)^n n^{n-1}}{n!}t^n
 =-t+t^2-\frac32t^3+\frac83t^4-\frac{125}{24}t^5
       +\frac{54}{5}t^6+\cdots.
\end{equation}
It is the infinitesimal branch $-W(t)$, but neither ordinary convergence of
partial sums in $\K$ nor a globally defined Lambert $W$ is required. Its
exact equation is $z_*+t e^{z_*}=0$.

\subsection{A double zero with nonstandard splitting}

For $F(z)=z^2-t$, the leading zero at $0$ has order two. The actual zeros
are $\pm t^{1/2}$, both infinitesimal and individually fine-isolated.
On an ordinary circle about zero,
\[
 \frac1{z^2-t}=\sum_{n\ge0}\frac{t^n}{z^{2n+2}}.
\]
Its coefficientwise contour integral is zero. The actual residues are
$1/(2t^{1/2})$ and $-1/(2t^{1/2})$, which are infinite but cancel exactly.
For $z/(z^2-t)$, the two actual residues are both $1/2$, so the integral is
$2\pi i$. For $F'/F=2z/(z^2-t)$, it is $4\pi i$, counting two zeros.

\subsection{Preparation of a genuinely analytic perturbation}

For $F(z)=z^2+t e^z$, the prepared polynomial has the form
$P=z^2+b_1(t)z+b_0(t)$. The recursion \eqref{eq:preprecursion} gives
\begin{align}
 b_0(t)&=t-\tfrac12t^2+\tfrac{11}{24}t^3+\bigo(t^4),\label{eq:prep-example0}\\
 b_1(t)&=t-\tfrac23t^2+\tfrac{27}{40}t^3+\bigo(t^4).\label{eq:prep-example1}
\end{align}
Here $\bigo(t^4)$ means a Hahn series supported in integer exponents at least
$4$, not a statement about limits of ordinary sequences. At first order,
\[
 p_1=1+z,\qquad v_1=\frac{e^z-1-z}{z^2}.
\]
At second order $H_2=-p_1v_1$, whose Taylor polynomial of degree below two is
$-1/2-(2/3)z$. At third order
$H_3=-(p_1v_2+p_2v_1)$, giving the displayed coefficients. The two roots of
$P$ lie in $\mm$ and are exactly the two zeros of $F$ there. Preparation has
reduced a transcendental local equation to a finite polynomial over $\K$.

\subsection{A high-rank perturbation}

Let
\[
 R(z)=t e^z+t^\omega\sin z,
 \qquad F(z)=z+R(z).
\]
The same simple-root formula applies. It automatically includes mixed
exponents from the monoid generated by $1$ and $\omega$. No assertion that
$n\to\infty$ eventually exceeds $\omega$ is made or needed. The
$t^\omega$ contributions cannot be justified by merely taking more terms
in a rank-one topological limit, but are handled by supportwise recursion.

\subsection{Scope of the computational checks}

The archive contains \texttt{verify\_examples.py}, which uses exact symbolic
arithmetic to verify the finite truncation of \eqref{eq:Wexample}, the first
three preparation coefficients in
\eqref{eq:prep-example0}--\eqref{eq:prep-example1}, the corresponding
factorization to the tested orders, and the rational residues in the
splitting example. These are reproducibility checks, not machine-checked
proofs of the general theorems. Infinite support arguments are supplied in
the mathematical proofs above.

\section{What the framework accomplishes, and what remains to develop}\label{sec:outlook}

The construction is useful because it makes a deliberate distinction between
three types of assertion. Algebraic closure and normal forms describe the
ambient numbers. Fine differentiation supplies local calculus. Coherent
holomorphic coefficients supply continuation and contour geometry. Classical
proofs can be transported only after identifying which of these structures
they use.

Preparation is particularly effective. It turns an arbitrary positive-support
analytic perturbation of an ordinary zero into a finite polynomial with
infinitesimal coefficients. This gives all nearby zeros, their exact total
multiplicity, rational principal parts of quotients, and an actual-pole
interpretation of the residue theorem. The argument never assumes ordinary
sequential completeness of $\K$.

The global exponential analysis gives a complementary lesson. The entire
unit circle is already reached by finite surreal angles. A global
conjugation-compatible exponential must therefore have additional periods
when infinite imaginary arguments are admitted. Local differential equations
and agreement on the ordinary complex plane do not determine those global
phases.

Several further directions are natural, but are not asserted as solved here.
One can develop locally supported sheaves and carefully specified atlases,
relating their cohomology to ordinary holomorphic coefficients. One can study
Runge approximation with explicit coefficient seminorms rather than the fine
topology. A theory of essential singularities would have to distinguish an
ordinary essential singularity in a coefficient from infinitely many moving
surcomplex poles. Finally, extra axioms on global phase maps, contour
families, and allowed growth scales could select a more rigid global
surcomplex analytic category. Each such choice must be checked against the
counterexamples already proved.

\appendix
\section{Dependency map and assumptions ledger}\label{app:ledger}

\begin{longtable}{@{}p{0.29\textwidth}p{0.64\textwidth}@{}}
\toprule
\textbf{Result} & \textbf{Essential input or restriction}\\
\midrule
\endhead
Taylor realization and calculus & A common well-ordered set support; complex
holomorphic coefficients; infinitesimal Taylor increments; Neumann's
positive-support lemma.\\[4pt]
Identity and maximum modulus & One connected ordinary coefficient domain;
not arbitrary locally analytic class functions.\\[4pt]
Preparation and zero lifting & A nonzero leading coefficient with a finite
ordinary zero order; positive-support perturbation; algebraic closure of
$\No[i]$.\\[4pt]
Inverse and open mapping & Perturbative inverse: a biholomorphic leading map.
General local mapping: infinitesimal rescaling and preparation; a nonzero
fine derivative gives the local inverse.\\[4pt]
Cauchy and Morera & Coefficientwise contours on ordinary curves; a common
support and ordinary coefficient regularity.\\[4pt]
Residues and removability & Quotients of Hahn-supported holomorphic functions.
For the contour theorem: leading denominator nonzero on the boundary of a
compactly contained bounded Jordan domain. Local boundedness removes poles
within this meromorphic category.\\[4pt]
Rouch\'e & Either a common leading coefficient or a uniform real strict
margin; zero counts occur in the full halo of the ordinary domain.\\[4pt]
Liouville & Boundedness of every ordinary coefficient, not merely one
surreal or real bound on the values of the whole function.\\[4pt]
Schwarz & The proved estimate is on standard parts; the exact classical
inequality fails on the specified halo disc.\\[4pt]
Strip exponential & Gonshor's real surreal exponential; standard complex
angles; infinitesimal formal exponential and logarithm.\\[4pt]
Global exponential obstruction & Group law, conjugation, and agreement with
the strip exponential; no analytic or topological hypothesis is needed.\\
\bottomrule
\end{longtable}

\noindent\textbf{Imported foundations.}
The real closure and normal-form structure of $\No$, the surreal real
exponential, Neumann's support lemma, the algebraic equivalence between real
closure and algebraic closure after adjoining $i$, and the ordinary complex
analysis theorems used coefficientwise are established inputs. All other
stated constructions and implications are proved in the article. This is not
a foundational reconstruction of surreal arithmetic, nor a formal proof
assistant development.

\section{Notation at a glance}
\begin{longtable}{@{}p{0.22\textwidth}p{0.71\textwidth}@{}}
\toprule
\textbf{Symbol} & \textbf{Meaning}\\
\midrule
\endhead
$\K=\No[i]$ & Surcomplex proper-class field.\\[3pt]
$t^\gamma=\omega^{-\gamma}$ & Conway normal-form monomial; exponent addition
is surreal addition.\\[3pt]
$\val(z)$ & Least exponent of a nonzero normal form.\\[3pt]
$z\prec w$ & $z/w$ is infinitesimal in modulus.\\[3pt]
$\OO,\mm$ & Finite surcomplex elements and their infinitesimal ideal.\\[3pt]
$\st$ & Standard-part homomorphism $\OO\to\C$.\\[3pt]
$U^\sharp$ & Halo $\st^{-1}(U)$ of an ordinary complex set.\\[3pt]
$\HH(U)$ & Globally Hahn-supported series with holomorphic coefficients on $U$.\\[3pt]
$D=\mathrm d/\mathrm dz$ & Derivative in an independent coordinate; every
surcomplex scalar coefficient is constant.\\[3pt]
$\cint$ & Coefficientwise Hahn contour integral over an ordinary shadow curve.\\[3pt]
$\Res_\zeta$ & Local Laurent residue at an actual surcomplex point.\\[3pt]
$\mathcal S$ & Strip $\No+i\OO_\R$ of finite imaginary part.\\[3pt]
$\Exp_\lambda$ & Explicit global exponential with phase projection
$p_\lambda(y)=y_{\fin}+\lambda c_\omega(y)$.\\
\bottomrule
\end{longtable}

\begin{thebibliography}{99}
\addcontentsline{toc}{section}{References}
\bibitem{Conway}
J. H. Conway, \emph{On Numbers and Games}, 2nd ed., A K Peters, 2001.
Foundational source for surreal arithmetic, normal forms, and surcomplex
numbers.

\bibitem{Gonshor}
H. Gonshor, \emph{An Introduction to the Theory of Surreal Numbers},
London Mathematical Society Lecture Note Series 110, Cambridge University
Press, 1986.

\bibitem{Alling}
N. L. Alling, \emph{Foundations of Analysis over Surreal Number Fields},
North-Holland Mathematics Studies 141, North-Holland, 1987.

\bibitem{Neumann}
B. H. Neumann, ``On ordered division rings'',
\emph{Transactions of the American Mathematical Society} \textbf{66}
(1949), 202--252.
\href{https://doi.org/10.1090/S0002-9947-1949-0032593-5}{doi:10.1090/S0002-9947-1949-0032593-5}.

\bibitem{vdH}
J. van der Hoeven, ``Operators on generalized power series'',
\emph{Illinois Journal of Mathematics} \textbf{45} (2001), 1161--1190.
\href{https://doi.org/10.1215/ijm/1258138061}{doi:10.1215/ijm/1258138061}.
Author version: \url{https://www.texmacs.org/joris/noeth/noeth.html}.

\bibitem{BMgerms}
A. Berarducci and V. Mantova, ``Transseries as germs of surreal functions'',
\emph{Transactions of the American Mathematical Society} \textbf{371}
(2019), 3549--3592.
\href{https://doi.org/10.1090/tran/7428}{doi:10.1090/tran/7428};
\href{https://arxiv.org/abs/1703.01995}{arXiv:1703.01995}.

\bibitem{BMderivations}
A. Berarducci and V. Mantova, ``Surreal numbers, derivations and transseries'',
\emph{Journal of the European Mathematical Society} \textbf{20} (2018),
339--390.
\href{https://doi.org/10.4171/JEMS/769}{doi:10.4171/JEMS/769};
\href{https://arxiv.org/abs/1503.00315}{arXiv:1503.00315}.

\bibitem{RSS}
S. Rubinstein-Salzedo and A. Swaminathan, ``Analysis on surreal numbers'',
\emph{Journal of Logic and Analysis} \textbf{6}:5 (2014), 1--39.
\href{https://arxiv.org/abs/1307.7392}{arXiv:1307.7392}.

\bibitem{CE}
O. Costin and P. Ehrlich, ``Integration on the surreals'',
\emph{Advances in Mathematics} \textbf{452} (2024), article 109823.
\href{https://arxiv.org/abs/2208.14331}{arXiv:2208.14331v5}.

\bibitem{MS}
J. M\"uller and A. Strohmaier, ``The theory of Hahn-meromorphic functions,
a holomorphic Fredholm theorem, and its applications'',
\emph{Analysis \& PDE} \textbf{7} (2014), 745--770.
\href{https://doi.org/10.2140/apde.2014.7.745}{doi:10.2140/apde.2014.7.745};
\href{https://arxiv.org/abs/1205.0236}{arXiv:1205.0236}.

\bibitem{Mantova2026}
V. Mantova, ``Monotonicity and a Taylor approximation theorem for
transseries'', January 2026,
\href{https://arxiv.org/abs/2601.07747}{arXiv:2601.07747}.

\bibitem{Ahlfors}
L. V. Ahlfors, \emph{Complex Analysis}, 3rd ed., McGraw--Hill, 1979.
Reference for the ordinary Cauchy, identity, Morera, residue, Rouch\'e,
maximum-modulus, Liouville, and Schwarz theorems used on coefficients.
\end{thebibliography}
\end{document}
